% Source: https://artofproblemsolving.com/community/c3h1411713_help_please
% Stable identity: aops-topic-1411713-q01
Let $V$ be the apex. The plane is $9-3=6$ units below $V$.

\begin{asy}
usepackage("amsmath");
size(5.5cm);
defaultpen(fontsize(10pt));
real a=6,h=9,t=3,q=0.35; real c=a*(h-t)/h; real bb=q*a;
pair V=(0,h),O=(0,0),P=(0,t);
real y0=bb^2/h, x0=a*sqrt(1-bb^2/h^2);
real th=aSin(bb/h);
draw(shift(O)*scale(a,bb)*arc((0,0),1,180-th,360+th),linewidth(0.9));
draw(shift(O)*scale(a,bb)*arc((0,0),1,th,180-th),dashed);
draw((-x0,y0)--V--(x0,y0),linewidth(0.9));
draw(shift(P)*scale(c,q*c)*arc((0,0),1,180-th,360+th),red+linewidth(0.9));
draw(shift(P)*scale(c,q*c)*arc((0,0),1,th,180-th),red+dashed);
draw(O--V,gray+dashed);
draw(O--(a,0),blue); draw(P--(c,t),blue);
dot(V);dot(O);dot(P);
label("$V$",V,N);
label("$6$",(a/2,0),N,blue);
label("$4$",(c/2,t),N,blue);
draw((-a-1.2,0)--(-a-1.2,t),Arrows(4)); label("$3$",(-a-1.2,t/2),W);
draw((-a-1.2,t)--(-a-1.2,h),Arrows(4)); label("$6$",(-a-1.2,(t+h)/2),W);
draw((-a-1.4,0)--(-a-1,0)); draw((-a-1.4,t)--(-a-1,t)); draw((-a-1.4,h)--(-a-1,h));
draw((-a-1,0)--(-a,0),gray+dotted); draw((-a-1,t)--(-c,t),gray+dotted); draw((-a-1,h)--V,gray+dotted);
label("\mbox{}",(min(currentpicture,user=true)+max(currentpicture,user=true))/2,black); // keeps the text after the figure black when LaTeX typesets the labels inline
\end{asy}

\textbf{(a)} Every point of the cone at distance $6$ below the apex lies at the same distance from the axis, so the cross-section is a \textbf{circle} (a disc) centred on the axis. The small cone above the plane is similar to the whole cone with ratio $\dfrac{6}{9}=\dfrac{2}{3}$, so the radius of the cross-section is $\dfrac{2}{3}\cdot6=4$ and its area is $\boxed{16\pi}$.

\textbf{(b)} The piece above the plane is a smaller \textbf{cone} (radius $4$, height $6$) and the piece below is a \textbf{frustum} (a truncated cone). Their volumes are
\[
\dfrac{1}{3}\pi\cdot4^2\cdot6=\boxed{32\pi}\qquad\text{and}\qquad \dfrac{1}{3}\pi\cdot6^2\cdot9-32\pi=108\pi-32\pi=\boxed{76\pi}.
\]
(Check with the scaling of volume: the small cone is $\left(\dfrac{2}{3}\right)^3=\dfrac{8}{27}$ of $108\pi$, which is $32\pi$.)

\textbf{(c)} A plane parallel to the base meets the cone (as a solid of finite height) only between the apex and the base. Strictly between them the section is always a circle; the plane of the base gives the base circle itself, and the plane through the apex meets the cone in a single \textbf{point} (a degenerate circle). So the possible figures are circles and, in the limiting case, a point.
